\chapter{Decidibilità, indecidibilità e parziale decidibilità}

\section{Predicati indecidibili}
Un predicato $M(\vec{x})$ si dice \textbf{decidibile} se la sua funzione caratteristica
$$
C_M(\vec{x}) = \begin{cases}
1 &\text{se $M(\vec{x})$ vale}\\
0 &\text{se $M(\vec{x})$ non vale} 
\end{cases}
$$
%
è calcolabile. Così non fosse, si direbbe indecidibile. Denominiamo l'algoritmo che calcola $c_M$ \textbf{procedura di decisione} del predicato $M(\vec{x})$. 

Altri modi per indicare che un predicato è decidibile, sono:
\it ricorsivo, ricorsivamente decidibile, risolvibile, ricorsivamente risolvibile, calcolabile, ha il problema della decisione ricorsivo\rm.

\section{Problemi indecidibili della teoria della calcolabilità}
\subsection{Appartenenza al dominio}
Il problema $x \in W_x$ è indecidibile, assieme a tutte le sue formulazioni equivalenti, come

$$
\phi_x(x)\downarrow,\quad P_x(x)\downarrow,\quad \psi_U(x,x)\downarrow
$$

\paragraph{Dimostrazione}
Dobbiamo dimostrare che la funzione

$$
f(x) = \begin{cases}
	1 &\text{se $x \in W_x$}\\
	0 &\text{se $x \notin W_x$} 
\end{cases}
$$
%
non è calcolabile. Procediamo per assurdo facendo vedere che se $f$ fosse calcolabile, riusciremo a costruire una $g$ diversa da tutte le funzioni calcolabili.

Supponiamo $f$ sia calcolabile, e definiamo

$$
g(x)
= \begin{cases}
	0 &\text{se $x \notin W_x$}\\
	\uparrow &\text{se $x \in W_x$}\\
\end{cases}
= \begin{cases}
	0 &\text{se $f(x) = 0$}\\
	\uparrow &\text{se $f(x) = 1$}\\
\end{cases}
$$

essendo una definizione per casi, questa deve essere calcolabile (imponendo che $f(x)$ sia calclabile, per ipotesi). Di contro, osserviamo che $x \in Dom(g) \Leftrightarrow x \notin W_x \ \forall x \in \mathbb{N}$, e quindi $Dom(g) \neq W_x \ \forall x \in \mathbb{N}$. $g$ è quindi diversa da tutte le funzioni calcolabili.

\subsection{Appartenenza al dominio e al range di \textit{h}}
Esiste una funzione calcolabile $h$ tale che i problemi $x \in Dom(h)$ e $x \in Ran(h)$ sono entrambi indecidibili.

\paragraph{Dimostrazione}

Si ponga
$$
g(x)
= \begin{cases}
	x &\text{se $x \in W_x$}\\
	\uparrow &\text{se $x \notin W_x$}\\
\end{cases}
$$
%
è una definizione per casi apparentemente calcolabile ($h(x) = x \cdot \underline 1 (\psi_U(x,x))$). Però, per ogni $x$ si ha 

$$
x \in Dom(h) \Leftrightarrow x \in W_x \Leftrightarrow x \in Ran(h)
$$

\subsection{Halting problem}
Il rinomato problema della fermata ``$\phi_x(y)\downarrow$'' è indecidibile. Lo dimostriamo usando una tecnica nota come \textbf{riduzione}. Informalmente, il problema $\phi_x(x)\downarrow$ è più facile del problema $\phi_x(y)\downarrow$. Poiché il primo è indecidibile, ne segue che anche il secondo lo è. Più formalmente, sia

$$
g(x,y)
= \begin{cases}
	x &\text{se $\phi_x(y)\downarrow$}\\
	\uparrow &\text{se $\phi_x(y)\uparrow$}\\
\end{cases}
$$
%
ora, se questa funzione $g$ fosse calcolabile, lo sarebbe anche il caso particolare $f(x) = g(x,x)$, ma questa sarebbe la funzione caratteristica del problema $\phi_x(x)\downarrow$ che è non calcolabile, pertanto $g$ non è calcolabile.

\subsubsection{Tecnica della riduzione}
Supponiamo di avere due problemi $M_1(\vec{x}),\, M_2(y_1, \dots, y_n)$. Diremo che $M_1$ è \textbf{riducibile} a $M_2$ se esistono funzioni totali calcolabili $k_1(\vec{x}), \dots, k_n(\vec{x})$ tali che 

$$
M_1(\vec{x}) \Leftrightarrow M_2(k_1(\vec{x}), \dots, k_n(\vec{x})) \ \forall \vec{x}
$$
%
Se $M_1$ è riducibile a $M_2$, ogni procedura di decisione per $M_2$ da luogo ad una procedura di decisione per $M_1$, per cui, se $M_1$ è indecidibile, anche $M_2$ dovrà essere indecidibile. Il teorema $s-m-n$ serve proprio a ridurre un problema a un altro.

\subsection{Indecidibilità di $\phi_x = \underline0$}
 $\phi_x = \underline0$ è indecidibile.
 
 \paragraph{Dimostrazione}
 Definiamo $f$ come 
 
 $$
 f(x,y) = \begin{cases}
 0 &\text{se $x \in W_x$}\\
 \uparrow &\text{se $x \notin W_x$}
 \end{cases}
 $$
 %
 Usiamo il teorema $s-m-n$ imponendo $x$ come prametro. $g_x(y) \simeq f(x,y)$. $g_x = \underline0 \Leftrightarrow x \in W_x$. Per la tesi di Church, $f$ è calcolabile, e quindi esiste una funzione calcolabile $k(x)$ tale che $f(x,y) \simeq \phi_{k(x)}(y)$, come $\phi_{k(x)} = g_x(y)$.
 
$$
x \in W_x \Leftrightarrow \phi_{k(x)} = \underline0
$$

\subsection{Indecidibilità di $\phi_x = \phi_y$}

\subsection{Indecidibilità dei problemi di input e di output}
Ossia, con $c \in \mathbb{N}$ $c \in W_x$ e $c \in E_x$

\section{Teorema di Rice}

Sia $\mathcal{B} \subseteq \mathcal{C}_1$, tale che $B \neq \emptyset$ e $\mathcal B \neq \mathcal C_1$. Il problema $\phi_x \in \mathcal B$ è indecidibile.
